\FloatBarrier
\section{Retrospective calibration audits added in revision v7}
\label{app:v7}
These analyses use existing predictions and already-consumed development data. They add no model fits, datasets, or external holdout openings. Their protocols and input/code hashes were recorded before the new computations. They are not prospective confirmation or population risk certification. The paired protocol specified the .95 relative cap and the full $2\times2$ head factorial; 220/220 is a reference configuration, not a separately prespecified primary cell.

\subsection{Experiment correspondence}
\begin{table}[H]\centering\scriptsize\setlength{\tabcolsep}{4pt}
\caption{Predictors, heads, caps and evidential roles. The main paired experiment has two predictors; the archived comparison has three different predictors. Reusing them in audits does not add independent model fits or replications. M5 dates below are dataset day indices.}\label{tab:experiment-map}
\begin{tabular}{p{.16\linewidth}p{.27\linewidth}p{.22\linewidth}p{.26\linewidth}}\toprule
Design & Predictor and head fitting & Calibration rule & Evaluation and role\\\midrule
Paired M5, two arms & Median LightGBM, 300 trees, days 1314--1554. Censored or complete histories/targets. Squared-loss error/exposure heads, 220/660 trees, days 1555--1610. & A 1611--1645; B 1646--1673. Primary $r=.95\max(R_A,R_B)$: .664708 censored; .632683 complete. Floor .35 in both. & Later 1800--1913. Seed 20260910; previously used items and periods. All eight head cells in Table~\ref{tab:matched-history}.\\
Archived M5, three predictors & XGBoost L1 hierarchy (750 trees, days 365--1313); hierarchy plus LightGBM L1 residual (selection 1314--1433, final depth 4/300 trees); frozen Chronos-2 median (context 2048, horizon 7, no task fine-tuning). Error/exposure heads: 220-tree squared-loss LightGBM, days 1555--1610. & Same A/B dates. Common reporting cap .640130; floor .35. Ordinary menu uses both windows. Archived screen designs on A at $.95r$ and checks 15 candidates on B. & Same Later window, supplementary predictor comparison; head seed 20260908. Separate from paired refits and the original external test (Appendix~\ref{app:external}).\\
V7 complete-sales screen & Reuses the complete-sales predictor and all four fitted head pairs. No refitting. & A design cap .601049; B screen at .632683; 20 candidates, tail .05/20. Utility uses A among B passers. & Previously used Later window; new conditional bootstrap summaries, no fresh confirmation.\\
V7 stability audit & Reuses complete-sales heads (four pairs) and the Bike models below. & Original caps fixed. Fixed thresholds: 4,000 draws. Threshold redesign: 200 draws for complete 220/220 and Bike. & Calibration only; no evaluation labels read. Item clusters for M5; day clusters for Bike.\\
Bike, supplementary & Median LightGBM, 200 trees, on 292 randomized days; Poisson exposure head on the same fit group; squared-loss error head on 146 other days. Seed 20260910. & 146 days (3,486 hours). Reporting cap .132806; threshold design .126166; floor .35. One calibration window. & 147 days (3,477 hours). Conditional count prediction using recorded weather, not temporal forecasting.\\\bottomrule
\end{tabular}\end{table}

\subsection{Calibration margins and selection stability}
\label{app:v7-stability}
The protocol at \texttt{evidence/revision\_v7\_stability/PROTOCOL.json} fixes both modes and their seeds before execution. All scores, fitted heads, reporting caps and $T$ stay fixed. Complete-sales calibration has 1,152,270 rows in 1,829 item clusters across A/B. One multinomial item multiplicity vector is shared by both windows, retaining all stores and days of an item. Bike resamples 146 calendar-day clusters and retains its original single-window design cap .95 times reporting risk. It preserves the original score scaling. No evaluation arrays are read.

Table~\ref{tab:v7-calibration} gives each candidate's original minimum case and exposure coverage and maximum risk. At complete 220/220, the case winner's margin over Mixed is 2.0109 points, while the exposure winner Ratio exceeds Descending by 0.02704 points. At the other complete-sales cells, Descending exceeds Ratio by 0.38843 (660/220), 0.09384 (220/660), and 0.36000 (660/660) points. These margins concern calibration utility, not evaluation performance or a perturbation bound for fitted moments.
\begin{table}[t]\centering\scriptsize
\caption{All five original calibration candidates in each complete-sales head cell and Bike. Coverage columns show the minimum over the original calibration windows (\%); Bike has one window. M5 reporting/design cap is .632683; Bike reporting cap is .132806 and design cap .126166. Error has no feasible threshold under the stated cap and .35 row floor; it remains in the menu record.}\label{tab:v7-calibration}
\begin{tabular}{llrrrl}\toprule
Trees $e/w$ & Rule & Min. case & Min. exposure & Max. risk & Feasible\\\midrule
220/220 & Error & 0.00 & 0.00 & -- & no \\
220/220 & Excess & 87.43 & 80.28 & 0.632681 & yes \\
220/220 & Mixed & 85.42 & 89.39 & 0.632683 & yes \\
220/220 & Ratio & 70.50 & 91.03 & 0.632683 & yes \\
220/220 & Descending & 48.67 & 91.00 & 0.632683 & yes \\
660/220 & Error & 0.00 & 0.00 & -- & no \\
660/220 & Excess & 86.85 & 79.10 & 0.632683 & yes \\
660/220 & Mixed & 85.90 & 88.49 & 0.632683 & yes \\
660/220 & Ratio & 75.36 & 90.61 & 0.632683 & yes \\
660/220 & Descending & 48.67 & 91.00 & 0.632683 & yes \\
220/660 & Error & 0.00 & 0.00 & -- & no \\
220/660 & Excess & 87.18 & 80.41 & 0.632682 & yes \\
220/660 & Mixed & 84.51 & 89.41 & 0.632683 & yes \\
220/660 & Ratio & 62.99 & 90.96 & 0.632683 & yes \\
220/660 & Descending & 48.67 & 91.05 & 0.632683 & yes \\
660/660 & Error & 0.00 & 0.00 & -- & no \\
660/660 & Excess & 87.31 & 79.65 & 0.632683 & yes \\
660/660 & Mixed & 85.64 & 88.80 & 0.632683 & yes \\
660/660 & Ratio & 69.56 & 90.69 & 0.632683 & yes \\
660/660 & Descending & 48.67 & 91.05 & 0.632683 & yes \\
Bike & Error & 0.00 & 0.00 & -- & no \\
Bike & Excess & 84.88 & 82.53 & 0.126117 & yes \\
Bike & Mixed & 83.99 & 88.10 & 0.126039 & yes \\
Bike & Ratio & 73.26 & 89.47 & 0.126156 & yes \\
Bike & Descending & 36.37 & 71.54 & 0.126142 & yes \\
\bottomrule\end{tabular}\end{table}

The four Descending case entries round to the same 48.67\%, but the two exposure heads do not accept the same rows. Their minimum window is A: 311,537/640,150 (48.666250\%) for 220 trees and 311,581/640,150 (48.673123\%) for 660. The A masks differ on 7,898 rows, with a net increase of 44 accepted rows; changing only the error head leaves Descending unchanged. The stricter A-only design uses different thresholds and accepts 224,362/224,653 rows (35.048348/35.093806\%). Independent row-level reconstruction confirms both tables; exact counts and hashes are in \texttt{docs/revision\_v9/DESCENDING\_COVERAGE\_CHECK.json}.

\paragraph{Fixed-threshold diagnostic.}
We retain each original threshold and use 4,000 cluster resamples (seed 20260913) to recheck positive accepted mass, empirical design risk, the .35 case floor, and both menu choices. M5 eligibility requires both original windows; Bike uses its single design window. At complete 220/220, Ratio/Descending exposure choices occur in 26.225\%/22.675\% of draws, and no-choice in 48.225\%. The central 95\% range of the signed Ratio-minus-Descending exposure margin, retaining both existing thresholds irrespective of eligibility, is [-0.4959,0.5297] points. Since thresholds were chosen at a calibration boundary, roughly half of fixed-threshold resamples need not remain eligible. This is not the failure rate of a procedure that is allowed to redesign its thresholds.

\paragraph{Threshold-redesign diagnostic.}
For complete 220/220 and Bike, 200 draws (seed 20260914) rerun the largest feasible whole-tie threshold for every ranking, then reselect both utilities. The implementation excludes score ties contributed only by zero-multiplicity clusters, reproduces all original thresholds when every multiplicity is one, and matches literal duplicated-row calculations on 40 toy cases including tied and infinite scores. This mode includes threshold and menu selection sensitivity conditional on the fixed fitted scores; it does not include fitting variability or cap estimation uncertainty.

Complete-sales exposure selection splits 103/97 between Ratio and Descending. Its signed exposure margin ranges from -0.4792 to 0.2810 points at the central 95\% empirical quantiles. Bike selects Ratio/Mixed in 192/8 exposure draws and Excess/Mixed in 161/39 case draws. Its signed Ratio-minus-Mixed exposure margin has central range [-0.0485,3.4899] points. No redesign draw has an empty menu. The 200 draws give at most about 3.54 percentage points of Monte Carlo standard error for a selection frequency; the results are coarse conditional sensitivity summaries. They do not establish expected deployment regret, population-optimal selection, or accepted-set overlap stability. The other three complete-sales head cells have only the fixed-threshold diagnostic, as specified before execution.
\begin{table}[t]\centering\scriptsize
\caption{Selection frequency (\%) with fixed original thresholds (4,000 cluster draws) or all thresholds redesigned within each resample (200 draws). Error, Excess, Mixed, Ratio, Descending and no-choice columns are all retained. Both modes hold the fitted scores, reporting cap and $T$ fixed. M5 clusters are items shared across A/B; Bike clusters are days. Fixed-mode failures are eligibility checks on boundary thresholds, not full-procedure failure rates.}\label{tab:v7-frequencies}
\begin{tabular}{lllrrrrrr}\toprule
Setting & Mode & Utility & Error & Excess & Mixed & Ratio & Desc. & None\\\midrule
220/220 & Fixed & c & 0.00 & 47.42 & 2.10 & 1.03 & 1.23 & 48.23 \\
220/220 & Fixed & d & 0.00 & 1.70 & 1.18 & 26.22 & 22.68 & 48.23 \\
220/220 & Redesign & c & 0.00 & 96.00 & 4.00 & 0.00 & 0.00 & 0.00 \\
220/220 & Redesign & d & 0.00 & 0.00 & 0.00 & 51.50 & 48.50 & 0.00 \\
660/220 & Fixed & c & 0.00 & 47.33 & 2.17 & 0.95 & 1.25 & 48.30 \\
660/220 & Fixed & d & 0.00 & 1.70 & 0.95 & 8.03 & 41.02 & 48.30 \\
220/660 & Fixed & c & 0.00 & 47.52 & 2.17 & 1.00 & 1.27 & 48.02 \\
220/660 & Fixed & d & 0.00 & 2.00 & 1.10 & 18.70 & 30.18 & 48.02 \\
660/660 & Fixed & c & 0.00 & 47.52 & 2.50 & 0.85 & 1.25 & 47.88 \\
660/660 & Fixed & d & 0.00 & 2.00 & 1.03 & 8.77 & 40.33 & 47.88 \\
Bike & Fixed & c & 0.00 & 44.62 & 10.27 & 1.77 & 6.05 & 37.28 \\
Bike & Fixed & d & 0.00 & 3.25 & 4.05 & 49.38 & 6.05 & 37.28 \\
Bike & Redesign & c & 0.00 & 80.50 & 19.50 & 0.00 & 0.00 & 0.00 \\
Bike & Redesign & d & 0.00 & 0.00 & 4.00 & 96.00 & 0.00 & 0.00 \\
\bottomrule\end{tabular}\end{table}

\begin{table}[t]\centering\small
\caption{Winner-minus-runner-up minimum-calibration-coverage margins (pp), among eligible candidates. The fixed-threshold interval is the empirical central 95\% resampling range conditional on at least two eligible candidates; $n$ gives that denominator out of 4,000. It is nonnegative by definition and does not measure whether a named pair switches order. Signed pair margins are reported in the text and machine-readable results.}\label{tab:v7-margins}
\begin{tabular}{llrlr}\toprule
Setting & Utility & Original margin & Fixed resampling range & $n$\\\midrule
220/220 & c & 2.01 & [1.48, 17.55] & 1944 \\
220/220 & d & 0.03 & [0.01, 2.34] & 1944 \\
660/220 & c & 0.95 & [0.46, 18.85] & 1943 \\
660/220 & d & 0.39 & [0.02, 3.15] & 1943 \\
220/660 & c & 2.67 & [2.08, 21.89] & 1935 \\
220/660 & d & 0.09 & [0.01, 1.94] & 1935 \\
660/660 & c & 1.67 & [1.22, 17.22] & 1940 \\
660/660 & d & 0.36 & [0.02, 2.51] & 1940 \\
Bike & c & 0.89 & [0.06, 44.94] & 2105 \\
Bike & d & 1.37 & [0.46, 12.74] & 2105 \\
\bottomrule\end{tabular}\end{table}

All candidate eligibility counts, both objectives' frequencies (including no-choice), unconditional and eligibility-conditional named-pair margins, and per-draw arrays are saved. Winner-minus-runner-up margins are nonnegative by construction; a signed named-pair margin is needed to describe an ordering switch. Evaluation outcomes are deliberately absent from this resampling audit.

\FloatBarrier
\subsection{Complete-sales A-design/B-screen audit}
\label{app:v7-screen}
The protocol at \texttt{evidence/revision\_v7\_complete\_screen/PROTOCOL.json} fixes $r=.6326834234117924$ from the existing complete-sales primary relative cap, score cap $r$, A design cap $.95r=.6010492522412028$, floor .35, four head pairs, and all five rankings. For each ranking, the largest eligible complete-tie threshold is designed on A alone. Thresholds are not retuned on B. At item $i$ on B, let
\[
 Z_i=\sum_{t:\,(i,t)\in B} a_{it} (L_{it}-r W_{it}).
\]
With 1,829 item clusters and 4,000 multinomial resamples (seed 20260912), $U^g_{20}$ is the .9975 empirical quantile of the resampled mean $Z_i$. An existing A candidate passes only if B accepted exposure is positive and $U^g_{20}\leq0$. The family size remains 20 including the four infeasible Error entries and duplicated Descending rules. Each utility selects maximum A coverage among passers, with the original fixed tie order and explicit no-choice fallback. FROZEN is written before EVALUATION\_STARTED and before reading Later arrays in this run. Hashing a previously available evaluation file is recorded separately from array access.

The full B outcome is in Table~\ref{tab:v7-B}. All 16 A-feasible candidates pass. Thus the B step removes none, and the stricter A cap plus changed design/selection procedure accounts for the changed policies; this audit does not identify a separate causal effect of B screening. The signed-excess criterion and ratio summaries use approximate cluster resampling, not the bounded independent-cluster theorem. A/B share items and common temporal shocks, and the reporting cap itself was previously chosen using A/B.
\begin{table}[t]\centering\scriptsize
\caption{Complete 20-candidate B screen. $U^{g}_{20}$ is the .9975 quantile of resampled mean signed excess per item, using reporting cap .632683 and A design cap .601049. Passing requires a defined A candidate, positive B exposure and $U^{g}_{20}\leq0$. All 16 A-feasible candidates pass; B removes no additional candidate. Error has undefined risk and is not eligible even though its empty excess is zero.}\label{tab:v7-B}
\begin{tabular}{llrrrrl}\toprule
Trees $e/w$ & Rule & A case (\%) & A exposure (\%) & B risk & $U^g_{20}$ & Pass\\\midrule
220/220 & Error & 0.00 & 0.00 & -- & 0.000 & no \\
220/220 & Excess & 75.18 & 70.47 & 0.59045 & -2.430 & yes \\
220/220 & Mixed & 68.72 & 80.07 & 0.59020 & -3.722 & yes \\
220/220 & Ratio & 50.65 & 82.94 & 0.59078 & -3.698 & yes \\
220/220 & Descending & 35.05 & 84.02 & 0.59063 & -4.013 & yes \\
660/220 & Error & 0.00 & 0.00 & -- & 0.000 & no \\
660/220 & Excess & 74.79 & 69.45 & 0.58963 & -2.524 & yes \\
660/220 & Mixed & 69.84 & 79.08 & 0.59067 & -3.443 & yes \\
660/220 & Ratio & 56.67 & 82.18 & 0.59097 & -3.520 & yes \\
660/220 & Descending & 35.05 & 84.02 & 0.59063 & -4.013 & yes \\
220/660 & Error & 0.00 & 0.00 & -- & 0.000 & no \\
220/660 & Excess & 74.80 & 70.73 & 0.59043 & -2.648 & yes \\
220/660 & Mixed & 66.92 & 80.36 & 0.59057 & -3.656 & yes \\
220/660 & Ratio & 45.06 & 82.86 & 0.59104 & -3.751 & yes \\
220/660 & Descending & 35.09 & 84.10 & 0.59052 & -4.046 & yes \\
660/660 & Error & 0.00 & 0.00 & -- & 0.000 & no \\
660/660 & Excess & 75.27 & 70.13 & 0.59086 & -2.477 & yes \\
660/660 & Mixed & 69.47 & 79.67 & 0.59122 & -3.498 & yes \\
660/660 & Ratio & 51.43 & 82.51 & 0.59052 & -3.783 & yes \\
660/660 & Descending & 35.09 & 84.10 & 0.59052 & -4.046 & yes \\
\bottomrule\end{tabular}\end{table}


On the previously consumed Later window, coverage and risk are evaluated for every candidate. Pointwise risk intervals use .025/.975 quantiles; one-sided upper summaries use .05/8 for the selected condition records and .05/20 for the full candidate family. The largest upper endpoints among the eight selected records are .616488 (pointwise), .621687 (eight-record adjustment), and .624450 (20-candidate adjustment), all below the fixed reporting cap. Duplicate Descending masks make these eight records six distinct policies; neither count represents independent replications. Table~\ref{tab:v7-cost} gives all selected records and costs, and Table~\ref{tab:v7-contrast} gives every contrast interval. Costs compare each cell's original utility-selected policy with its new A-design/B-screen choice, so the reference exposure cost includes a change from Ratio to Descending as well as a threshold change.
\begin{table}[t]\centering\scriptsize
\caption{All eight selected policies on Later. Case utility selects Excess; exposure utility selects Descending. Costs are unscreened minus A-design/B-screen coverage (pp), relative to each cell\textquotesingle s original utility-selected policy. Risk intervals and $U_8$ (tail .05/8) are approximate and fit-conditional. These shared-head policies are not eight independent replications.}\label{tab:v7-cost}
\begin{tabular}{llrrrlrrr}\toprule
Trees & Utility & Case \% & Exp. \% & Risk & Pointwise 95\% CI & $U_8$ & Cost $c$ & Cost $d$\\\midrule
220/220 & c & 73.55 & 70.81 & 0.59428 & [0.57538, 0.61566] & 0.62056 & 12.77 & 10.14 \\
220/220 & d & 37.16 & 83.89 & 0.59787 & [0.58072, 0.61641] & 0.62084 & 33.34 & 6.78 \\
660/220 & c & 72.99 & 69.60 & 0.59446 & [0.57533, 0.61572] & 0.62088 & 12.68 & 9.99 \\
660/220 & d & 37.16 & 83.89 & 0.59787 & [0.58072, 0.61641] & 0.62084 & 13.99 & 7.14 \\
220/660 & c & 73.27 & 71.19 & 0.59484 & [0.57611, 0.61567] & 0.62078 & 12.87 & 10.04 \\
220/660 & d & 37.25 & 83.96 & 0.59794 & [0.58082, 0.61649] & 0.62087 & 14.04 & 7.15 \\
660/660 & c & 73.55 & 70.43 & 0.59550 & [0.57655, 0.61645] & 0.62169 & 12.73 & 9.97 \\
660/660 & d & 37.25 & 83.96 & 0.59794 & [0.58082, 0.61649] & 0.62087 & 14.04 & 7.15 \\
\bottomrule\end{tabular}\end{table}

\begin{table}[t]\centering\small
\caption{A-design/B-screen Later contrasts and descriptive central 95\% paired item-bootstrap intervals, in percentage points. These compare the exposure-selected and case-selected policies after the full procedure.}\label{tab:v7-contrast}
\begin{tabular}{lrlrl}\toprule
Trees $e/w$ & $\Delta c$ & Interval & $\Delta d$ & Interval\\\midrule
220/220 & -36.39 & [-38.77, -34.02] & +13.08 & [10.90, 15.50] \\
660/220 & -35.83 & [-38.22, -33.42] & +14.29 & [11.88, 17.12] \\
220/660 & -36.02 & [-38.42, -33.62] & +12.77 & [10.72, 15.00] \\
660/660 & -36.29 & [-38.68, -33.87] & +13.53 & [11.14, 16.35] \\
\bottomrule\end{tabular}\end{table}

These calculations supply simultaneous-family-adjusted \emph{approximate summaries}, not a demonstrated simultaneous coverage theorem. They condition on fitted models and selected thresholds, reuse development outcomes, and omit common store/calendar dependence beyond item clustering. They do not validate the original unscreened policy pair or its union. No new union risk interval was computed in this audit. No head tuning, additional training seed, learned SelectiveNet baseline, new non-WAPE positive example, or prospective holdout was added.

Implementation repair details, retained earlier attempts, raw-archive recovery records, and independent verification reports are provided in the accompanying source package under \texttt{docs/revision\_v7/} and the audit evidence directories. They are not additional scientific trials.
